%%%%%%%%%%%%%%%%%%%%%%%%%%%%%%%%%%%%%%%%
% Here's some basic formatting packages:
%%%%%%%%%%%%%%%%%%%%%%%%%%%%%%%%%%%%%%%%
\usepackage{graphicx}                       % Required for inserting images
\usepackage{xcolor}                         % some lovely color additions!!
\setlength{\parskip}{0.5em}                  % this forces LaTeX to make paragraph skips not be single-spaced
%\usepackage{import}                         % this allows us to bring separate .tex files into the main document using \import, \input, etc.
%\usepackage{lscape}                         % this allows us to put things into landscape via \begin{landscape}
\usepackage{setspace}                       % this allows us to put block quotes into single spaced typeset using the following:
\usepackage[normalem]{ulem}                 % this allows for strikethrough text using \sout{}
\AtBeginEnvironment{quote}{\par\singlespacing\small} % handling big block quotes

%colors!
\definecolor{UTorange}{RGB}{255,130,0} % Define Tennessee Orange
\definecolor{smokeygray}{RGB}{75,75,75} % Define Smokey Gray
\definecolor{UTMlimestone}{RGB}{240,237,227} % Define UT Martin Limestone

\usepackage{hyperref}                       % this allows us to color links using the following:
\hypersetup{
    colorlinks=true,
    citecolor=black,
    linkcolor=black,
    filecolor=UTorange,      
    urlcolor=UTorange,
    }
\urlstyle{same}
\usepackage[margin=1in]{geometry}  % this makes the size of the page consistent and sets the document-wide margin sizes


%%%%%%%%%%%%%%%%%%%%%%%%%%%%%%%%%%%%%%%%%%%%%%%%%
% Here's the packages that manage my bibliography
%%%%%%%%%%%%%%%%%%%%%%%%%%%%%%%%%%%%%%%%%%%%%%%%%
\usepackage{xurl}
\usepackage[style=apa,backend=biber,url=true]{biblatex}
\addbibresource{references.bib} % Replace with your .bib file

% Limit citations to 2 authors before using et al.
\DeclareNameFormat{author}{%
  \ifnum\value{listcount}=1
    \ifnum\value{liststop}>2
      \usebibmacro{name:last}{##1}{##3}{##5}{##7}%
      \space et~al.%
    \else
      \usebibmacro{name:last-first}{##1}{##3}{##5}{##7}%
    \fi
  \else
    \usebibmacro{name:last-first}{##1}{##3}{##5}{##7}%
  \fi
  }


% use full author field on first cite, then shortauthor on following cites
\makeatletter
\AtEveryCitekey{%
  \ifciteseen
    {}
   {\clearfield{shortauthor}}%
}
\makeatother
  
  % this is a footnote without a number
  \newcommand\blfootnote[1]{%
  \begingroup
  \renewcommand\thefootnote{}\footnote{#1}%
  \addtocounter{footnote}{-1}%
  \endgroup
}


%%%%%%%%%%%%%%%%%%%%%%%%%%%%%%%%%%%%%%%%%%%%%%%%
% Here's the packages for good math typesetting:
%%%%%%%%%%%%%%%%%%%%%%%%%%%%%%%%%%%%%%%%%%%%%%%%
\usepackage{amsmath}                    % this enables all the nice math stuff, like \frac and \begin{bmatrix}
\usepackage{amsthm}                     % this enables theory work like follows:
\newtheorem{prop}{Proposition}
\newtheorem{cor}{Corollary}[prop]
\usepackage{amsfonts}                   % this enables the fancy expectation operator as \mathbb{E}, among other symbols
%\usepackage{bbm}                        % this allows lowercase blackboard bold as \mathbbm{y}, and so on
%\usepackage{bm}                         % enables generalized bolded math-mode symbols via \bm{text}
\allowdisplaybreaks                     % this command allows \begin{align} environments to extend over multiple pages
\DeclareMathOperator*{\argmax}{argmax}  % this gives the nice under-argmin operator via \argmin_{thing under argmin}
\DeclareMathOperator*{\argmin}{argmin}  % this gives the nice under-argmin operator via \argmin_{thing under argmin}

%%%%%%%%%%%%%%%%%%%%%%%%%%%%%%%%%%%%%%%%%%%%%%%%
% Here's the packages for nice tables & figures:
%%%%%%%%%%%%%%%%%%%%%%%%%%%%%%%%%%%%%%%%%%%%%%%%
\usepackage{float}                          % this allows positioning of tables via [h!], etc.
\usepackage{threeparttable}                 % this makes for easy, consistently-formatted notes under tables that I love and appreciate dearly
\usepackage{booktabs}                       % this enables topline, cmidrule, etc. in our tables
\usepackage{arydshln}                       % this enables \hdashline
\usepackage{caption}                        % this properly spaces table captions an appropriate distance from the tables
\usepackage{subcaption}                     % this gives us good notes beneath figures using \subcaption or \subfigure
\usepackage{adjustbox}                      % this allows us to put nice boxes around figures and tables using \adjustbox{} if we want to do so
